\clearpage
\section*{Preliminary Work}

\subsection*{Problem 1. Full-Wave Bridge Rectifier}


\textbf{a)} Positive half cycle: D1 and D3 \\
Negative half cycle: D2 and D4

\textbf{b)} 100 Hz because the Full wave rectifier folds the negative voltage to positive so the VL frequency is twice the input.

\textbf{c)} Peak voltage is 9 volts (Vs=9) so according to $2V_{sp}/\pi$ formula the average voltage is 5.730v but because of the diodes which are in series the total voltage drops by 1.2 volts so the average voltage drops to 4.530 volts. 
The the DC load current can be calculated with $V_{LDC}/R_L$ which is equal to 3.775mA. 
In the end the total power will be calculated as $2 \cdot V_D \cdot I_{LDC} = 2 \cdot 0.6\text{ V} \cdot 3.775\text{mA} = 4.53\text{mW}$.

\subsection*{Problem 2. Full-wave Rectifier with a Center-tapped Transformer}
\begin{enumerate}
    \item Average DC Load Voltage: $V_{LDC} = 18/\pi - 0.6 = 5.130\text{V}$
    \item Average DC Load Current: $I_{LDC} = 5.130/1200 = 4.275\text{ mA}$
    \item Total Power: $1 \cdot V_D \cdot I_{LDC} = 0.6 \cdot 4.275 = 2.57\text{ mW}$
\end{enumerate}

\subsection*{Problem 3. Full-wave Bridge Rectifier with Triangular Approximation}
\textbf{a) Load Regulation ($R_L = 330\text{ ohms}$)}
\begin{enumerate}
    \item \textbf{No-load condition ($V_{NL}$):} When R is infinity no current is drawn so the capacitor does not charge and voltage remains at peak. \\
    $V_{NL} = V_{LP} = 7.8\text{V}$
    \item \textbf{Full Load Condition:} $f_r R_L C = 100 \cdot 330 \cdot 22 \cdot 10^{-6} = 0.726$ \\
    $V_{FL} = 7.8 / (1 + (1 / (2 \cdot 0.726))) = 4.619\text{V}$
    \item \textbf{Calculate Load Regulation:}
    \item $\text{Load Reg} = \left(\frac{V_{NL} - V_{FL}}{V_{NL}}\right) 100\% = \left(\frac{7.8 - 4.619}{7.8}\right) \cdot 100\% = 40.78\%$
\end{enumerate}

\textbf{b) Calculate Ripple Factor}
\begin{itemize}
    \item \textbf{For Load Resistance = 330 ohms:} \\
    Ripple Factor = $1 / (3.4641 \cdot 100 \cdot 330 \cdot 0.000022)$ \\
    Ripple Factor = 0.3976 \\
    Result = 39.76\%
    
    \item \textbf{For Load Resistance = 1.2 kohms (1200 ohms):} \\
    Ripple Factor = $1 / (3.4641 \cdot 100 \cdot 1200 \cdot 0.000022)$ \\
    Ripple Factor = 0.1093 \\
    Result = 10.93\%
    
    \item \textbf{For Load Resistance = 10 kohms (10000 ohms):} \\
    Ripple Factor = $1 / (3.4641 \cdot 100 \cdot 10000 \cdot 0.000022)$ \\
    Ripple Factor = 0.0131 \\
    Result = 1.31\%
\end{itemize}

\textbf{c) Effect of Load Resistance on Ripple Voltage:} \\
As the load resistance increases, the load draws less current from the circuit. A lower current draw means the filter capacitor discharges much slower between the active charging cycles. Because the capacitor loses less stored voltage before the next pulse arrives, the peak-to-peak voltage drop (the ripple) is significantly reduced, resulting in a cleaner, smoother DC output.

